\textbf{Program 1: All three probing methods in one program.}

\textbf{File:} \texttt{hashing\_collisions.cpp}.

\textbf{Chapter connection:} Linear probing; Quadratic probing; Double hashing; Searching under open addressing.

Share the probe rule between insertion and search. Start with a fresh table for each method and count every inspected slot. The prime table sizes and nonnegative keys keep the double-hash step valid; insertion stops after at most m probes.

\begin{adslabfigure}[How this program runs]
\begin{adslabgraphic}
\begin{tikzpicture}[font=\small,execute at begin node={\hyphenpenalty=10000\exhyphenpenalty=10000}]
\node[utpstep,text width=2.50cm,minimum height=1.15cm] (s0) at (0.0,0) {\textbf{1}\\Choose probe rule};
\node[utpstep,text width=2.50cm,minimum height=1.15cm] (s1) at (2.95,0) {\textbf{2}\\Insert bounded probes};
\node[utpstep,text width=2.50cm,minimum height=1.15cm] (s2) at (5.9,0) {\textbf{3}\\Search same sequence};
\node[utpstep,text width=2.50cm,minimum height=1.15cm] (s3) at (8.85,0) {\textbf{4}\\Compare tables / costs};
\draw[utparrow] (s0.east)--(s1.west);
\draw[utparrow] (s1.east)--(s2.west);
\draw[utparrow] (s2.east)--(s3.west);
\end{tikzpicture}
\end{adslabgraphic}
\end{adslabfigure}


\begin{utpsteps}\setcounter{enumi}{0}\item \textbf{Headers and initial definitions.}\end{utpsteps}

\texttt{hashing\_collisions.cpp} --- lines 18--34.

\lstinputlisting[firstline=18,lastline=34]{all_experiments/Lab11/hashing_collisions.cpp}


\tcbbreak

\begin{utpsteps}\setcounter{enumi}{1}\item \textbf{Implement probeIndex, insertKey.}\end{utpsteps}

\texttt{hashing\_collisions.cpp} --- lines 35--63.

% Worked lab input; chapter fonts, colours and TikZ styles retained.
\begin{adslabfigure}[Insert 15 by double hashing: home 0, step 4]
\begin{adslabgraphic}
\begin{tikzpicture}[font=\small]
\node[utphot,minimum width=0.88cm,minimum height=.65cm] at (0.00,0.00) {10};
\node[utplabel] at (0.00,0.60) {0};
\node[utpdim,minimum width=0.88cm,minimum height=.65cm] at (0.92,0.00) {--};
\node[utplabel] at (0.92,0.60) {1};
\node[utpcell,minimum width=0.88cm,minimum height=.65cm] at (1.84,0.00) {12};
\node[utplabel] at (1.84,0.60) {2};
\node[utpdim,minimum width=0.88cm,minimum height=.65cm] at (2.76,0.00) {--};
\node[utplabel] at (2.76,0.60) {3};
\node[utpcell,minimum width=0.88cm,minimum height=.65cm] at (3.68,0.00) {14};
\node[utplabel] at (3.68,0.60) {4};
\node[utplabel,anchor=west,align=left,text width=5.2cm] at (4.8,0.00) {i=0: $(0+0\times4)\bmod5=0$\\Occupied by 10: continue};
\draw[utparrow] (-.85,-0.35)--(-.85,-0.85);
\node[utpcell,minimum width=0.88cm,minimum height=.65cm] at (0.00,-1.20) {10};
\node[utpdim,minimum width=0.88cm,minimum height=.65cm] at (0.92,-1.20) {--};
\node[utpcell,minimum width=0.88cm,minimum height=.65cm] at (1.84,-1.20) {12};
\node[utpdim,minimum width=0.88cm,minimum height=.65cm] at (2.76,-1.20) {--};
\node[utphot,minimum width=0.88cm,minimum height=.65cm] at (3.68,-1.20) {14};
\node[utplabel,anchor=west,align=left,text width=5.2cm] at (4.8,-1.20) {i=1: $(0+1\times4)\bmod5=4$\\Occupied by 14: continue};
\draw[utparrow] (-.85,-1.55)--(-.85,-2.05);
\node[utpcell,minimum width=0.88cm,minimum height=.65cm] at (0.00,-2.40) {10};
\node[utpdim,minimum width=0.88cm,minimum height=.65cm] at (0.92,-2.40) {--};
\node[utpcell,minimum width=0.88cm,minimum height=.65cm] at (1.84,-2.40) {12};
\node[utphot,minimum width=0.88cm,minimum height=.65cm] at (2.76,-2.40) {15};
\node[utpcell,minimum width=0.88cm,minimum height=.65cm] at (3.68,-2.40) {14};
\node[utplabel,anchor=west,align=left,text width=5.2cm] at (4.8,-2.40) {i=2: $(0+2\times4)\bmod5=3$\\Empty: store 15 and return 3};
\end{tikzpicture}
\end{adslabgraphic}
\end{adslabfigure}


\lstinputlisting[firstline=35,lastline=63]{all_experiments/Lab11/hashing_collisions.cpp}


\tcbbreak

\begin{utpsteps}\setcounter{enumi}{2}\item \textbf{Implement findKey, list, printTable and helpers.}\end{utpsteps}

\texttt{hashing\_collisions.cpp} --- lines 64--92.

\lstinputlisting[firstline=64,lastline=92]{all_experiments/Lab11/hashing_collisions.cpp}


\tcbbreak

\begin{utpsteps}\setcounter{enumi}{3}\item \textbf{Connect the test cases in main.}\end{utpsteps}

\texttt{hashing\_collisions.cpp} --- lines 93--126.

\textbf{Read the chapter visual with the code.} The diagram illustrates the storage rule; the program tests the values and cases stated in \texttt{main}.
% Copied from chapters/11-hashing.tex (fig:probing).
% Keep the chapter drawing, type sizes, colours, and node styles unchanged.
\begin{utpcompactlabfigure}[Linear probing and clustering]
\begin{labgraphic}
\begin{tikzpicture}[font=\small]
\foreach \row/\hot/\lbl [count=\r from 0] in {%
  {10,{},12,{},14}/{9}/{10, 12, 14 land directly: 1 probe each},
  {10,15,12,{},14}/{1}/{$15\bmod5=0$ taken $\rightarrow$ probe 1: free},
  {10,15,12,20,14}/{3}/{$20\bmod5=0$ taken $\rightarrow$ probes 1, 2, 3}}{
  \begin{scope}[yshift=-1.35*\r cm]
    \foreach \v [count=\i from 0] in \row{
      \node[utplabel] at (1.2*\i,0.58) {\i};
      \ifnum\i=\hot
        \node[utphot,minimum width=1.12cm,minimum height=0.72cm] at (1.2*\i,0) {\v};
      \else
        \ifx\v\empty
          \node[utpdim,minimum width=1.12cm,minimum height=0.72cm] at (1.2*\i,0) {};
        \else
          \node[utpcell,minimum width=1.12cm,minimum height=0.72cm] at (1.2*\i,0) {\v};
        \fi
      \fi
    }
    \node[utplabel,anchor=west,align=left] at (5.4,0) {\lbl};
  \end{scope}
}
% cluster bracket on last row
\draw[decorate,decoration={brace,mirror,amplitude=4pt,raise=3pt},draw=UTPRed,semithick]
  (-0.62,-3.11) -- (4.42,-3.11);
\node[utplabel,anchor=north,color=UTPRed] at (1.9,-3.55)
  {a cluster: no free slot until the table wraps};
\end{tikzpicture}
\end{labgraphic}
\end{utpcompactlabfigure}


\textbf{Trace checkpoint for the supplied main.} After inserting 10, 12 and 14 into a size-five table, slots 0, 2 and 4 are occupied. For key 15, double hashing uses step 4 and visits 0, 4, then 3. Search must reproduce that same sequence. In the crowded-table exercise, distinguish failure to reach a free slot within the probe bound from a table that has no free slot at all.

\lstinputlisting[firstline=93,lastline=126]{all_experiments/Lab11/hashing_collisions.cpp}


\tcbbreak

\begin{utpsteps}\setcounter{enumi}{4}\item \textbf{Compare probe counts in a crowded table.}\end{utpsteps}

\texttt{hashing\_collisions.cpp} --- lines 127--152. Continue in the same file; do not add another header or opening brace.



\lstinputlisting[firstline=127,lastline=152]{all_experiments/Lab11/hashing_collisions.cpp}


\tcbbreak

\textbf{Build, run, and inspect hashing\_collisions.}

From this program\textquotesingle s folder, run these commands in PowerShell:

\begin{lstlisting}[style=utpterminal,language={}]

g++ -std=c++17 -Wall -Wextra hashing_collisions.cpp -o hashing_collisions.exe

.\hashing_collisions.exe

\end{lstlisting}

\textbf{Expected output excerpt.} Compare values, order, and pass/fail messages; spacing need not match.

\begin{lstlisting}[style=utpterminal,language={}]

== H(x) = x mod 5, keys 10 12 14 15 ==

linear probing:
  key 10  home 0  probes 0               -> slot 0
  key 12  home 2  probes 2               -> slot 2
  key 14  home 4  probes 4               -> slot 4
  key 15  home 0  probes 0 -> 1          -> slot 1
  table: [0:10] [1:15] [2:12] [3:-] [4:14] 
  [ok]   every key is found again with the same probe rule
  [ok]   every probe index stayed inside the table

quadratic probing:
  key 10  home 0  probes 0               -> slot 0
  key 12  home 2  probes 2               -> slot 2
  key 14  home 4  probes 4               -> slot 4

\end{lstlisting}

The complete captured run is in \texttt{labs/outputs/hashing\_collisions.txt}.

\begin{tryit}

Use a fresh size-5 table with quadratic inserts 0,5,10,15. The final route 0,1,4,4,1 fails despite free slots. Record failure separately from the five probes and confirm the table is unchanged.

\end{tryit}
